% Geometric Brownian Motion — research-note PDF.
%
% PROSE AND CODE ARE VERBATIM from the website article:
%   site/src/content/articles/brownian-motion.tsx   (prose, bullets, captions)
%   site/src/content/articles/gbm-spy-code.ts       (all code cells + stdout)
% Nothing is paraphrased, condensed or added.
%
% Metadata (title, dek, project card) from site/src/lib/articles.ts.
% Figures are Louis's, carried over from site/public/figures/.
%
% LAYOUT follows the reference note's own convention, measured from it:
% section heads sit full width at the left margin (x=39.6 in every case) and
% usually open a page; the body runs in two 255.4bp columns beneath. Code,
% display equations and figures break out to the full measure because they do
% not fit a column — the longest code line needs 452bp against 231bp of column.
%
% Each section is its own two-column block, so a head can never strand at the
% foot of a column belonging to the previous section, and \noteneedroom keeps a
% head from landing without room for real content under it.
%
% NOTE: the article has no disclosures section, so this note carries none.
\documentclass{pyportfolios-note}

\noteshorttitle{Geometric Brownian Motion}

\begin{document}
\thispagestyle{notecover}

% ============================================================ cover ==========
% Kicker is the registry's format field verbatim. An earlier draft read
% "Tutorial · Quant Foundations", abbreviating cardCategory ("Quant Foundations
% & Derivatives") into a form that appears nowhere in the source.
% Q3 2026 is derived from the registry date, 2026-08-12.
%
% The second right-hand line is deliberately empty. The reference carries the
% data period there, but it has no project card; ours does, and its ASSETS and
% TIMEFRAME rows already state "SPY (S&P 500 ETF, State Street)" and
% "Jan 2018 – Dec 2024" on this same page. Repeating them in the masthead said
% the same thing twice, 450bp apart.
\notemasthead{Quantitative Insights}{Tutorial}{Q3 2026}{}

\notetitle{Geometric\\Brownian Motion}

\notedeck{Simulating price paths with the continuous-time workhorse behind
Black–Scholes and Monte-Carlo risk engines.}

\textbf{Key Highlights:}

\begin{multicols}{2}
\begin{notebullets}
\item \textbf{Geometric Brownian Motion (GBM)} is the canonical continuous-time
      model for asset prices.
\item GBM powers the \textbf{Black–Scholes} framework and \textbf{Monte-Carlo}
      risk engines.
\item It entails most of the \textbf{intuition practitioners} carry about
      ``what could the price do?''.
\item The framework serves as fundamental bedrock of mathematical finance —
      advanced applications layer additional features on top of it.
\end{notebullets}
\end{multicols}

% <ProjectCard slug="brownian-motion" /> — sits under the highlights, full
% width, as the article places it. Five rows, no format row, per the component.
\notecard
  {Quant Foundations \& Derivatives}
  {NumPy · Pandas · Matplotlib · yfinance · SciPy}
  {SPY (S\&P 500 ETF, State Street)}
  {Jan 2018 – Dec 2024}
  {Scenario cones for wealth projections; the engine inside Monte-Carlo pricing
   \& risk systems}

% The cover ends here by design: opening a 26-line setup listing on it would
% leave a four-line stub under the card and split the cell across two pages.
\newpage

% ============================================================ setup ==========
\begin{notecode}
# Setup

import numpy as np
import pandas as pd
import matplotlib.pyplot as plt

np.random.seed(42)  # reproducibility (standard arbitrary number)
plt.rcParams["figure.dpi"] = 110

# ---- house style ---------------------------------------------------------
BG    = "#F4F2E8"   # Pale Sisal
INK   = "#151515"   # Anthracite
AQUA  = "#1FFFFF"   # accent
GREY  = "#7A7A7A"

HEAD  = {"fontname": "Switzer", "fontsize": 11, "fontweight": "bold", "color": INK}

plt.rcParams.update({
    "font.family": "monospace",
    "font.monospace": ["Courier Prime", "Courier New", "DejaVu Sans Mono"],
    "font.size": 8,
    "axes.titlesize": 11, "axes.labelsize": 8,
    "xtick.labelsize": 7.5, "ytick.labelsize": 7.5, "legend.fontsize": 7.5,
    "figure.facecolor": BG, "axes.facecolor": BG, "savefig.facecolor": BG,
    "axes.edgecolor": GREY, "axes.labelcolor": GREY,
    "xtick.color": GREY, "ytick.color": GREY,
    "axes.spines.top": False, "axes.spines.right": False,
    "axes.grid": True, "grid.alpha": 0.3, "grid.color": "#CCCCCC", "grid.linewidth": 0.5,
    "legend.frameon": False,
    "lines.linewidth": 2.0, "lines.markersize": 0,
})
\end{notecode}

% ======================================================== 1 Introduction =====
\noteneedroom{150bp}
\notesection{1.\notenumsep Introduction}

A few simple facts to get up to speed:

\begin{multicols}{2}
\begin{notebullets}
\item GBM assumes that \textbf{percentage returns} — not price changes — are
      random, normally distributed, and independent over time.
\item Prices therefore stay positive and their terminal distribution is
      \textbf{log-normal}.
\item \textbf{2 parameters} fully describe the model — both estimable directly
      from historical data:
      \begin{notesubbullets}
      \item the drift $\mu$ (expected annual log-growth plus half-variance)
      \item the volatility $\sigma$ (annualised standard deviation of log returns)
      \end{notesubbullets}
\end{notebullets}
\end{multicols}

% ========================================================== 2 Intuition ======
% Room for the head, its lead line and first bullet group — not the whole
% section. Demanding the whole section (260bp) forces it to the next page and
% costs a 190bp hole here, a 280bp hole later and a 38%-full final page; the
% section instead runs across the page break, which is ordinary in two columns.
\noteneedroom{140bp}
\notesection{2.\notenumsep Intuition}

The logic before diving into equations:

\begin{multicols}{2}
\begin{notebullets}
\item \textbf{Returns compound, prices don't add}
      \begin{notesubbullets}
      \item A \$10 move means something very different at SPY = 100 than at
            SPY = 600.
      \item What is comparable across price levels is the \textbf{relative} move.
      \item GBM makes randomness proportional to the current price:
            $dS = \mu S\,dt + \sigma S\,dW$.
      \end{notesubbullets}
\item \textbf{Why log-normal}
      \begin{notesubbullets}
      \item If each day multiplies the price by a small random gross return, the
            log-price is a \textbf{sum} of many small independent shocks.
      \item Sums of independent shocks are (approximately) normal.
      \item Normal log-price $\Rightarrow$ log-normal price: skewed right,
            floored at zero.
      \end{notesubbullets}
\item \textbf{Drift vs. noise}
      \begin{notesubbullets}
      \item Over short horizons the noise term dominates: $\sigma\sqrt{t}$
            shrinks slower than $\mu t$ as $t \to 0$.
      \item Over long horizons drift wins.
      \item This is the same $\sqrt{t}$ rule that governs plain Brownian
            motion — GBM simply wraps it in an exponential.
      \end{notesubbullets}
\end{notebullets}
\end{multicols}

% ============================================================= 3 Theory ======
% Full measure throughout: the display equations run to ~320bp and would not
% set in a 231bp column.
\noteneedroom{200bp}
\notesection{3.\notenumsep Theory}

Below the theoretical foundation and respective equations:

\begin{notebullets}
\item \textbf{The Stochastic Differential Equation (SDE) and its solution}
      \begin{notesubbullets}
      \item A differential equation with a random term — so it yields a
            \textit{distribution} of paths, not only one.
      \item $\mu S_t\,dt$ is the \textbf{drift} (deterministic trend);
            $\sigma S_t\,dW_t$ is the \textbf{diffusion} (the noise).
      \end{notesubbullets}
      \[ dS_t = \mu S_t\,dt + \sigma S_t\,dW_t \]
\item \textbf{Applying Itô's lemma to $\ln S_t$}
      \begin{notesubbullets}
      \item The $-\tfrac{1}{2}\sigma^2$ correction comes from $(dW)^2 = dt$.
      \end{notesubbullets}
      \[ S_t = S_0 \exp\!\Big[\big(\mu - \tfrac{1}{2}\sigma^2\big)t + \sigma W_t\Big] \]
\item \textbf{Exact discretisation}
      \begin{notesubbullets}
      \item Because the solution is closed-form, we can simulate
            \textbf{without discretisation error} at any step size $\Delta t$.
      \end{notesubbullets}
      \[ S_{t+\Delta t} = S_t \cdot
         \exp\!\Big[\big(\mu - \tfrac{1}{2}\sigma^2\big)\Delta t
         + \sigma \sqrt{\Delta t}\, Z\Big], \qquad Z \sim \mathcal{N}(0,1) \]
\end{notebullets}

The key implementations of those small functions: one draws per-step
\textbf{gross returns}, the other \textbf{compounds} them into price paths.

\begin{notecode}
# Implementation

def gbm_returns(mu, sigma, dt, n_steps, n_paths):
    """One-step gross returns under GBM (shape: n_steps x n_paths)."""
    z = np.random.normal(size=(n_steps, n_paths))
    return np.exp((mu - 0.5 * sigma**2) * dt + sigma * np.sqrt(dt) * z)

def gbm_paths(s0, mu, sigma, dt, n_steps, n_paths):
    """Price paths of shape (n_steps + 1, n_paths), starting at s0."""
    rets = gbm_returns(mu, sigma, dt, n_steps, n_paths)
    return s0 * np.vstack([np.ones(rets.shape[1]), rets]).cumprod(axis=0)
\end{notecode}

% ======================================================== 4 Application ======
% Full measure: alternating listings, printed output and full-width figures.
\noteneedroom{150bp}
\notesection{4.\notenumsep Application}

\begin{notebullets}
\item \textbf{Calibrate $\mu$ and $\sigma$ from real data}
      \begin{notesubbullets}
      \item We pull daily SPY prices, estimate annualised drift and volatility
            from \textbf{log returns}, and let the data drive the simulation.
      \end{notesubbullets}
\end{notebullets}

\begin{notecode}
TICKER = "SPY"
START, END = "2018-01-01", "2024-12-31"

import yfinance as yf

px = (yf.download(TICKER, start=START, end=END, auto_adjust=True, progress=False)["Close"]
        .squeeze().rename(TICKER).dropna())

log_ret = np.log(px / px.shift(1)).dropna()

s0        = float(px.iloc[-1])
sigma_hat = float(log_ret.std() * np.sqrt(252))
mu_hat    = float(log_ret.mean() * 252 + 0.5 * sigma_hat**2)  # drift of the SDE, not of log returns

print(f"{TICKER}: {len(px)} daily observations, last close = {s0:,.2f}")
print(f"mu_hat    = {mu_hat:.2%}  (annualised drift)")
print(f"sigma_hat = {sigma_hat:.2%}  (annualised volatility)")
\end{notecode}

\begin{notecodeout}
SPY: 1760 daily observations, last close = 578.32
mu_hat    = 14.75%  (annualised drift)
sigma_hat = 19.54%  (annualised volatility)
\end{notecodeout}

\noteneedroom{120bp}
\begin{notebullets}
\item \textbf{Simulate 1,000 5-year paths}
      \begin{notesubbullets}
      \item Grey lines are individual paths, the dashed line is the mean of the
            simulated distribution, and the band spans the 5th–95th percentile.
      \end{notesubbullets}
\end{notebullets}

\begin{notecode}
DT      = 1 / 252
HORIZON = 252 * 5          # 5 years
N_PATHS = 1_000

paths = gbm_paths(s0, mu_hat, sigma_hat, DT, HORIZON, N_PATHS)
t_ax  = np.arange(paths.shape[0]) / 252

p5, p95 = np.percentile(paths, [5, 95], axis=1)

fig, ax = plt.subplots(figsize=(10, 5))
ax.fill_between(t_ax, p5, p95, color=AQUA, alpha=0.35, lw=0, zorder=1,
                label="5\u201395% band")
ax.plot(t_ax, paths[:, :200], color="#BBBBBB", lw=0.35, alpha=0.45, zorder=2)
ax.plot(t_ax, paths.mean(axis=1), color=INK, lw=2.2, ls="--", zorder=3,
        label="Mean path")
ax.set_title(f"{TICKER}: {N_PATHS:,} GBM paths, 5 years  "
             f"(mu={mu_hat:.1%}, sigma={sigma_hat:.1%})", **HEAD)
ax.set_xlabel("Years"); ax.set_ylabel("Simulated price"); ax.legend()
plt.tight_layout(); plt.show()
\end{notecode}

% Chart generated by figures/make_brownian_motion.py — the article's own
% gbm_paths() under its seed and printed parameters, styled to the reference
% note's chart palette. Caption text is the article's, with the figure number
% and the source line added per the reference's figure convention.
\notefigure{figures/brownian-motion/fig1-paths.png}
  {Figure 1. SPY: 1,000 GBM paths, 5 years (mu=14.7\%, sigma=19.5\%)}
  {Simulated price under geometric Brownian motion, daily steps; the band spans
   the 5th–95th percentile of 1,000 paths}
  {Source: Author calculations; simulated, not actual, results. 1,000 paths from
   the article's calibrated parameters ($S_0$ = 578.32, $\mu$ = 14.75\%,
   $\sigma$ = 19.54\%), estimated on SPY daily closes, January 2018 – December 2024.}

\noteneedroom{170bp}
\begin{notebullets}
\item \textbf{Terminal distribution} — is it log-normal?
      \begin{notesubbullets}
      \item \textbf{Verifies the code:} the simulated histogram should match the
            closed-form density
            $\ln\!\big(S_T/S_0\big) \sim \mathcal{N}\big((\mu - \tfrac{1}{2}\sigma^2)T,\; \sigma^2 T\big)$
            — if it doesn't, the ½σ² term is usually the culprit.
      \item \textbf{Proves the claim:} section 2 asserts prices end up
            log-normal; this is where you see the right skew rather than take it
            on faith.
      \item \textbf{The practical lesson:} the \textbf{mean sits above the
            median}, so the ``average'' projected outcome is not the
            \textit{typical} one — most paths land below it.
      \end{notesubbullets}
\end{notebullets}

\begin{notecode}
from scipy import stats

T   = HORIZON / 252
s_T = paths[-1]

x   = np.linspace(s_T.min(), s_T.max(), 400)
pdf = stats.lognorm.pdf(x, s=sigma_hat*np.sqrt(T),
                        scale=s0*np.exp((mu_hat - 0.5*sigma_hat**2)*T))

fig, ax = plt.subplots(figsize=(10, 4))
ax.hist(s_T, bins=60, density=True, color=AQUA, alpha=0.55,
        edgecolor=INK, linewidth=0.25, label="Simulated $S_T$")
ax.plot(x, pdf, color=INK, lw=2.2, label="Theoretical log-normal")
ax.axvline(s_T.mean(),     color=INK, lw=1.8, ls="--", label=f"Mean   {s_T.mean():,.0f}")
ax.axvline(np.median(s_T), color=INK, lw=1.8,          label=f"Median {np.median(s_T):,.0f}")
ax.set_title(f"{TICKER}: terminal price distribution after {T:.0f} years", **HEAD)
ax.set_xlabel("Price"); ax.legend()
plt.tight_layout(); plt.show()
\end{notecode}

\notefigure{figures/brownian-motion/fig2-terminal.png}
  {Figure 2. SPY: terminal price distribution after 5 years}
  {Simulated terminal prices against the closed-form log-normal density; the
   mean sits above the median, the signature of the right skew}
  {Source: Author calculations; simulated, not actual, results. Distribution of
   $S_T$ across the same 1,000 paths, with the closed-form log-normal density
   overlaid.}

\noteneedroom{190bp}
\begin{notebullets}
\item \textbf{Sanity check:} switch the drift off
      \begin{notesubbullets}
      \item With $\mu = 0$ the process becomes a \textbf{martingale} — no
            expected gain — so the mean terminal price must land back at $S_0$.
      \item Zeroing one parameter isolates the rest: this catches a missing
            $-\tfrac{1}{2}\sigma^2$ correction that a full simulation might hide.
      \item \textbf{How to read the result:} the sample mean wobbles around
            $S_0$ with standard error
            $\approx S_0\sqrt{e^{\sigma^2T}-1}\big/\sqrt{N}$. Within \textasciitilde2 SE
            is noise; a persistent gap in one direction that \textit{doesn't}
            shrink as $N$ grows is a bug.
      \end{notesubbullets}
\end{notebullets}

\begin{notecode}
paths_nd = gbm_paths(s0, 0.0, sigma_hat, DT, HORIZON, N_PATHS)
print(f"Mean terminal price (no drift): {paths_nd[-1].mean():,.2f}   vs   S0 = {s0:,.2f}")
\end{notecode}

\begin{notecodeout}
Mean terminal price (no drift): 588.06   vs   S0 = 578.32
\end{notecodeout}

% ========================================================= 5 Conclusion ======
% The four groups read as one unit (~290bp); move them together rather than
% tailing a single group onto a page of its own.
% This one genuinely must not split: at 140bp the block began on the previous
% page and stranded "Alternatives & Extensions" — head on one page, its three
% bullets on the next at 4% fill. ~270bp moves all four groups together.
\noteneedroom{270bp}
\notesection{5.\notenumsep Conclusion}

\begingroup
% Four short labelled groups; at the document's 7bp \parskip the paragraph
% gaps alone cost ~30bp and push the block onto a page of its own at 40% fill.
% 3bp still reads as spacious in a 255bp column.
\setlength{\parskip}{3bp}
\begin{multicols}{2}
\textbf{Strengths}

\begin{notebullets}
\item \textbf{Analytically tractable} — closed-form solutions (Black–Scholes)
      make it the natural baseline
\item \textbf{Positive prices, log-normal terminal distribution} — matches the
      basic stylised fact that prices can't go negative
\item \textbf{Only two parameters} — both estimable directly from historical log
      returns
\item \textbf{Cheap to simulate} — vectorised NumPy generates millions of paths
      in seconds
\end{notebullets}

\textbf{Weaknesses \& Limitations}

\begin{notebullets}
\item \textbf{Constant volatility} — realilty has vol clusters and spikes
\item \textbf{No jumps} — crashes like March 2020 are far outside its reach
\item \textbf{Normal log-returns} — real returns have fat tails
\item \textbf{Independent increments} — momentum and mean-reversion exist
\end{notebullets}

\textbf{Applications in Practice}

\begin{notebullets}
\item Baseline for option pricing and Monte-Carlo risk engines
\item Scenario cones for wealth projections
\item Base model against which fancier models must justify their complexity
\end{notebullets}

\textbf{Alternatives \& Extensions}

\begin{notebullets}
\item \textbf{Stochastic volatility} — Heston
\item \textbf{Jump-diffusion} — Merton
\item \textbf{GARCH-family models} — for clustered volatility
\end{notebullets}
\end{multicols}
\endgroup

\end{document}
